\documentclass{gescons}

\genre {Entrevista}
\author{Denise Paro}
\title{Autoconscientização Grafopensênica}


\begin{document}
    \makeentrevistatitle
    \coverart{back/Denise_Paro}

    \begin{multicols}{2}

\begin{center}
    \includegraphics[width=8cm]{articles/entrevista/mockups/Denise_Paro.png}
\end{center}

\subsubsection{1. Qual foi a motivação para a escrita da obra? Por que a definição deste tema para publicação de um livro?}

A motivação surgiu a partir de uma parapercepção em uma atividade da
Dinâmica Parapsíquica da Seriexologia, da Consecutivus, instituição
dedicada às pesquisas de vidas passadas. Acessei a seguinte frase:
``Escrever para fazer justiça'' e percebi um grupo de escritores
desorientados e sentindo-se injustiçados porque escreviam para defender
causas sociais e humanitárias, mas acabaram tolhidos e incompreendidos.

Tais personalidades não entendiam o motivo pelo qual eram atacadas se
tinham a intenção de ajudar pessoas e apoiar causas.

A partir desse parafato, uma série de sincronicidades ocorreram que
culminaram na decisão de escrever o livro a fim de evidenciar como a
escrita pode ser ferramenta assistencial, porém sem a necessidade de
\emph{carregar nas tintas} e \emph{atacar diretamente as pessoas}. A
crítica pode ser feita, porém de modo cosmoético e assistencial.

\subsubsection{2. Quais foram as principais percepções, intra e extrafísicas, durante a pesquisa e a escrita da obra? E posterior ao lançamento?}

Ficou evidente o acoplamento com amparadores que me inspiraram com
ideias para alguns tópicos contidos no sumário do livro e nas
argumentações. Uma série de sincronicidades ocorreram no sentido de
facilitar a imersão no tema, boa parte delas relacionadas ao Caso
Dreyfus e ao contexto dos escritores engajados que perpassa a tese do
livro.

\begin{pullquote}
``Uma série de sincronicidades ocorreram no sentido de facilitar a imersão no tema.''
\end{pullquote}

Após o lançamento, ficou ainda mais claro meu retropapel na escrita
engajada e como isso ainda se manifesta nesta atual vida intrafísica,
porém no atual momento com mais lucidez no sentido da
autoconscientização quanto aos efeitos do que escrevo.

\begin{pullquote}
``Após o lançamento, ficou ainda mais claro meu retropapel na escrita engajada e como isso ainda se manifesta nesta atual vida intrafísica.''
\end{pullquote}

\subsubsection{3. Qual o maior aprendizado com a escrita desta obra?}

A lucidez quanto a um dos meus retrogrupos do passado, no caso de
escritores engajados, e a necessidade de se fazer uma ponte assistencial
para atender tal público. Daí o título da obra Autoconscientização
Grafopensênica, no sentido de que é importante estar cientes dos efeitos
do que se escreve e como um texto ou discurso, quer que seja, pode
contribuir para formar holopensenes mais hígidos ou patológicos.

\subsubsection{4. O que poderia dizer como incentivo para que mais pesquisadores invistam na publicação de obras conscienciológicas?}

Um dos principais efeitos da obra conscienciológica é o da
autoconscientização da condição atual da conscin nesta vida intrafísica
quanto ao tema estudado e a oportunidade de se fazer reciclagens
conscienciais. Portanto, uma oportunidade evolutiva ímpar.

    
    \end{multicols}
\end{document}
